% SPDX-License-Identifier: CC-BY-4.0
% Copyright (c) 2026 Martin Schröder <info@swedishembedded.com>
%
% The page-1 figure: the S4 procedure in five steps beside its result, drawn from the ledgered
% data (generated/rsi4_pct.csv, generated/rsi3_pct.csv). Set in the title area, not floated.

\begin{center}
\begin{minipage}{\textwidth}
\centering
{\large\bfseries Retrained from the base on a doubling frontier, every generation beats the last in 4 of 4
lineages (registered per-item test) on an exam the learner never sees.\par}
\vspace{2mm}
\begin{tikzpicture}[font=\normalsize, >={Stealth[length=2.6mm]}]
\tikzset{hstep/.style={draw=black!75, rounded corners=3pt, fill=#1, align=center, text width=40mm,
    minimum height=11mm, inner sep=2.5pt, line width=0.6pt, font=\small},
  hstep/.default=okSky!25,
  hnum/.style={circle, fill=black!80, text=white, font=\footnotesize\bfseries, inner sep=1.2pt,
    minimum size=4.6mm}}
\node[hstep] (p) at (0,0) {\textbf{Model proposes}\\Lua programs\\from the manual};
\node[hstep=okGreen!18] (l) at (0,-15mm) {\textbf{Interpreter labels}\\each program};
\node[hstep] (f) at (0,-30mm) {\textbf{Keep the frontier}:\\wrong in $\ge$1 of 4 tries};
\node[hstep=okPurple!22] (t) at (0,-45mm) {\textbf{Fresh adapter from}\\\textbf{the base} on the pool\\(300 $\to$ 900 $\to$ 2100)};
\node[hstep=okOrange!25] (e) at (0,-61mm) {\textbf{Separate examiner}\\grades a held-out exam};
\foreach \a/\b in {p/l, l/f, f/t, t/e} \draw[->, line width=0.9pt, black!75] (\a) -- (\b);
\draw[->, line width=0.9pt, dashed, black!55] (t.east) -- ++(5mm,0) |- node[pos=0.25, right, font=\footnotesize, align=left] {next\\round} (p.east);
\foreach \n/\k in {p/1, l/2, f/3, t/4, e/5} \node[hnum, anchor=east] at ([xshift=-1.5mm]\n.west) {\k};
\begin{axis}[paper, at={(60mm,-67mm)}, anchor=south west, width=87mm, height=80mm,
  xmin=-0.2, xmax=3.45, ymin=0, ymax=30, xtick={0,1,2,3}, xticklabels={base,1,2,3},
  xlabel={generation}, ylabel={held-out items right (\% of 889)},
  tick label style={font=\small}, label style={font=\small},]
\addplot[okBlue, line width=1.6pt, mark=*, mark size=2.2pt] table[col sep=comma,x=gen,y=r0] {generated/rsi4_pct.csv};
\addplot[okBlue, line width=1.6pt, mark=*, mark size=2.2pt, forget plot] table[col sep=comma,x=gen,y=r1] {generated/rsi4_pct.csv};
\addplot[okBlue, line width=1.6pt, mark=*, mark size=2.2pt, forget plot] table[col sep=comma,x=gen,y=r2] {generated/rsi4_pct.csv};
\addplot[okBlue, line width=1.6pt, mark=*, mark size=2.2pt, forget plot] table[col sep=comma,x=gen,y=r3] {generated/rsi4_pct.csv};
\addplot[okVermillion, line width=1.1pt, densely dashed, mark=o, mark size=2pt] table[col sep=comma,x=gen,y=r0] {generated/rsi3_pct.csv};
\addplot[okVermillion, line width=1.1pt, densely dashed, mark=o, mark size=2pt, forget plot] table[col sep=comma,x=gen,y=r1] {generated/rsi3_pct.csv};
\addplot[okVermillion, line width=1.1pt, densely dashed, mark=o, mark size=2pt, forget plot] table[col sep=comma,x=gen,y=r2] {generated/rsi3_pct.csv};
\addplot[okVermillion, line width=1.1pt, densely dashed, mark=o, mark size=2pt, forget plot] table[col sep=comma,x=gen,y=r3] {generated/rsi3_pct.csv};
\node[anchor=west, font=\normalsize\bfseries, text=okBlue] at (axis cs:3.07,24.3) {S4};
\node[anchor=west, font=\normalsize\bfseries, text=okVermillion] at (axis cs:3.07,13.2) {S3};
\end{axis}
\end{tikzpicture}
\captionof{figure}{\textbf{S4 at a glance.} Left: one round; rounds share no parameters, only the growing
pool. Right: share of 889 held-out items answered (items the base had missed at selection) by the four S4 lineages (solid), retrained from the base each round, and the four S3 lineages (dashed), which continued one adapter. Under
intervals clustered over the exam's 25 templates the third step holds in one or two lineages, and a
registered matched-data control, not yet run, decides whether the gain needs the improved proposer or
only more verified data (\cref{sec:rsi,sec:followups}).}\label{fig:hero}
\end{minipage}
\end{center}
\vspace{1mm}
